\documentclass{ar2rc}

% ==============================
% Packages & bibliography
% ==============================
\usepackage{commands}
\usepackage{silence}

\BiblatexSplitbibDefernumbersWarningOff
\WarningFilter{latex}{Empty bibliography}
\addbibresource{ref.bib}

% ==============================
% Link to the manuscript used for jumps + appendix
% ==============================
\setmainpaper{demo/main}

% ==============================
% Metadata
% ==============================
\papertitle{Optimizing Neural Networks for Edge Devices: A Comprehensive Study}
\manusid{TPAMI-2025-01-1234}
\authors{San Zhang$^{1}$, Si Li$^{1}$, Wu Wang$^{2}$, and Liu Zhao$^{2}$}
\affiliations{%
$^{1}$ Department of Computer Science, University of Technology, City, Country.\\
$^{2}$ AI Research Lab, Innovation Corp, City, Country.%
}
\contactemails{san.zhang@univ.edu, si.li@univ.edu, wu.wang@innocorp.com, liu.zhao@innocorp.com}

\journalname{IEEE Transactions on Pattern Analysis and Machine Intelligence}
\editorname{Prof. John John}
\resubmissiondate{\today}
\decisiondate{January 15, 2026}

\begin{document}
  %===============================================================================
  %                               Cover Letter
  %===============================================================================

  \responseletterheader

  \begin{flushleft}
    \printresubmissiondate\\
    \printeditorname\\
    Editor-in-Chief\\
    \printjournalname\\[0.8em]
    Dear \printeditorname,
  \end{flushleft}

  \vspace{0.8em}

  Thank you for your email of \printdecisiondate, regarding our manuscript entitled
  ``\textit{\printpapertitle},'' originally submitted to the \textit{\printjournalname}.
  We have carefully revised the manuscript in response to the comments provided by
  the Associate Editor and the three reviewers.
  Based on your suggestion, we are now submitting the revised manuscript as a
  \textit{new submission}.
  For your reference, the original log number is \printmanusid.

  For clarity and ease of reading, the following formatting conventions are used
  throughout this response:

  \begin{itemize}
    \item Reviewers' comments are presented in \textit{italic font} to distinguish
      them from the responses.
    \item The responses appear in regular font.
    \item Revised sentences in the manuscript with tracked changes are highlighted
      in {\color{blue} blue}.
    \item Blue badges such as \revisionbadge{r1c1a} are clickable and jump between
      this letter and the appended manuscript.
    \item When revised text is quoted in this response letter, it is shown in
      {\color{blue} blue} for consistency.
  \end{itemize}

  We look forward to hearing from you and sincerely appreciate your time and
  consideration.

  \vspace{1.5em}

  \responselettersignature
  \newpage

\section{Responses to the Comments from Editor}

\begin{generalblock}
\begin{generalcomment}
  \label{com0.0}
  Please ensure that all figures are of high quality (at least 300 dpi) and that
  the references follow the journal's style guide.
\end{generalcomment}
\begin{generalresponse}
  We have double-checked all figures and replaced low-resolution images with
  vector graphics or high-resolution bitmaps where appropriate. We have also
  verified the reference formatting to ensure compliance with the journal's
  guidelines.
\end{generalresponse}
\end{generalblock}

\section{Responses to the Comments from Reviewer 1}
\resetcommentcounter

\begin{generalblock}
\begin{generalcomment}
  \label{com1.0}
  The paper presents an interesting approach to optimizing neural networks for
  edge devices. However, the experimental validation needs to be strengthened,
  particularly regarding the comparison with state-of-the-art methods.
\end{generalcomment}
\begin{generalresponse}
  Thank you for your positive feedback and constructive suggestions. We have
  significantly expanded the experimental section by including comparisons with
  three additional state-of-the-art methods~\cite{EtAlForever2020}.
\end{generalresponse}
\end{generalblock}

\begin{reviewerblock}
\begin{reviewercomment}
  \label{com1.1}
  It is unclear how the proposed method handles varying network architectures.
  Could the authors provide more details on the generalization capability of
  their approach?
\end{reviewercomment}
\begin{reviewerresponse}
  We appreciate this insightful comment. Our method is designed to be
  architecture-agnostic. We revised the manuscript in \rlinks{r1c1a,r1c1b}.

  \revisionblock{Our approach is architecture-agnostic: the same optimization
  pipeline applies to CNNs, Transformers, and hybrid models without hand-tuned
  architecture-specific rules.}
\end{reviewerresponse}
\end{reviewerblock}

\begin{reviewerblock}
\begin{reviewercomment}
  \label{com1.2}
  Figure 3 is difficult to read. The font size of the axis labels is too small.
\end{reviewercomment}
\begin{reviewerresponse}
  We apologize for the oversight. We have redrawn Figure~3 with larger fonts.
  The updated figure is marked at \rlink{r1c2}.
\end{reviewerresponse}
\end{reviewerblock}

\newpage
\section{Responses to the Comments from Reviewer 2}
\resetcommentcounter

\begin{generalblock}
\begin{generalcomment}
\label{com2.0}
This paper addresses a timely problem. The proposed method is technically sound,
but the introduction could be improved to better motivate the research.
\end{generalcomment}
\begin{generalresponse}
Thank you for your encouraging comments. We have rewritten the Introduction to
better highlight the challenges in deploying large neural networks on
resource-constrained edge devices and to clearly state our contributions.
\end{generalresponse}
\end{generalblock}

\begin{reviewerblock}
\begin{reviewercomment}
\label{com2.1}
In Equation (5), the variable $\lambda$ is introduced but not defined.
\end{reviewercomment}
\begin{reviewerresponse}
Thank you for pointing this out. We defined $\lambda$ immediately after
Equation~(5). See \rlink{r2c1}.
\end{reviewerresponse}
\end{reviewerblock}

\begin{reviewerblock}
\begin{reviewercomment}
\label{com2.2}
The ablation study in Table 2 is informative, but it would be better to include
the inference time as a metric.
\end{reviewercomment}
\begin{reviewerresponse}
We agree with the reviewer. We have updated Table~2 to include the average
inference time (in milliseconds) measured on a Raspberry Pi 4 device. The
results show that our method achieves a 20\% reduction in latency compared to
the baseline.
\end{reviewerresponse}
\end{reviewerblock}

\newpage
\section{Responses to the Comments from Reviewer 3}
\resetcommentcounter

\begin{generalblock}
\begin{generalcomment}
\label{com3.0}
The paper is well-written and easy to follow. I have only a few minor comments.
\end{generalcomment}
\begin{generalresponse}
We thank the reviewer for the positive assessment and the careful reading of our
manuscript.
\end{generalresponse}
\end{generalblock}

\begin{reviewerblock}
\begin{reviewercomment}
\label{com3.1}
Page 4, line 15: ``network are'' should be ``network is''.
\end{reviewercomment}
\begin{reviewerresponse}
Fixed. Thank you for catching this typo. See \rlink{r3c1}.
\end{reviewerresponse}
\end{reviewerblock}

\begin{reviewerblock}
\begin{reviewercomment}
\label{com3.2}
Please check the formatting of Reference~\cite{EtAlForever2020}.
\end{reviewercomment}
\begin{reviewerresponse}
We have corrected the formatting of Reference~\cite{EtAlForever2020} to ensure
it matches the journal's citation style.
\end{reviewerresponse}
\end{reviewerblock}

\appendmainpaper

\end{document}
